\documentclass[phy1200]{handout}

\title{Diffraction of light}

\begin{document}

\maketitle

Light being a wave instead of particles leads to some surprising behavior. The main such behavior is the \textbf{interference} of waves, also known as diffraction. Overlapping waves that are in phase (peaks line up with peaks) or out of phase (peaks line up with troughs) can add together or cancel each other out, respectively. 

For an overview of interference and double slit diffraction, see the textbook\footnote{Open Stax university physics volume 3 chapter 3.1-2: \href{https://openstax.org/books/university-physics-volume-3/pages/3-1-youngs-double-slit-interference}{Link to textbook}}. The important formula is the spacing between bright spots far from the slits:
$$
\sin(\theta) = \frac{n\lambda}{d},
$$
where $\theta$ is the angle of the ray relative to the beam direction, $\lambda$ is the wavelength of the light, and $d$ is distance between the slits.

In this demonstration, you will use a laser and some diffraction slits to see the relationship between slit spacing, wavelength, and the diffraction pattern.

\section{Known diffraction slits}
You will be doing diffraction with a laser.\mgpar{Be safe: be careful where you shine the laser. You'll be shining it across the classroom, so make sure nobody is in the path of the beam so you don't shine it in their eye.}

Shine the laser at your paper from a distance of about 10 cm and note the size of the beam. 

Shine the laser at a wall across the room and note the how the size of the beam has changed.

Use the diffraction slit apparatus with a red and green laser. Now trace the beam pattern, noting the bright and/or dark spots. Try a few slit spacings (there is even a part where you can continuously vary the slit spacings), and see how these change the diffraction pattern. How does the color (wavelength) of the laser change the pattern? 

\section{Unknown diffraction slits}
The important part of diffraction isn't that the slits are slits, just that there is a blockage between the two openings. As a result, you can accomplish the same with a single obstacle, such as a hair. 

Pluck a hair from your head and hold it in the laser beam. You should see a diffraction pattern. 

You can actually measure the distance between dark spots to figure out the diameter of the hair. Consider how to convert the spot spacing into a value of $\sin\theta$. 

\section{Report}
In your report, be sure to address the following:
\begin{enumerate}
    \item If light were made of particles, what would be the image you would expect as light passes through two slits?
    \item Explain interference, and why the two beams have a phase difference. 
    \item Explain why $d\sin\theta$ has to do with the phase difference.
    \item How does diffraction spacing change when you change the slit spacing?
    \item How does diffraction spacing change when you change the wavelength of the light?
    \item Explain your method of figuring out the width of a human hair and your result.
\end{enumerate}


\end{document}